\documentclass[10pt,a4paper]{amsart}
\usepackage[margin=19mm]{geometry}
\usepackage{amsmath,amssymb,mathtools}
\usepackage[scr=rsfs,cal=euler]{mathalfa}
\usepackage{xfrac}
\usepackage[unicode,hidelinks,bookmarks=false]{hyperref}
\newcommand{\ie}{i.e.\ }
\newcommand{\cf}{cf.\ }
\newcommand{\resp}{respectively\ }
\newcommand{\Subsection}{\subsection}
\providecommand{\nobreakdash}{\nobreak\hbox{-}}
\setlength{\emergencystretch}{2em}
\multlinegap0pt
\usepackage{bm}
\usepackage{bbm}
\usepackage{dsfont}
\usepackage{xspace}
\usepackage{cancel}
\usepackage{mathtools}
\usepackage{amsthm}
\makeatletter
\@ifundefined{theorem}{\newtheorem{theorem}{Theorem}}{}
\@ifundefined{lemma}{\newtheorem{lemma}[theorem]{Lemma}}{}
\makeatother
\newcommand{\op}[1]{\operatorname{#1}}
\title[On the number of subrings of $\mathbb{Z}^n$ of prime power i]{On the number of subrings of $\mathbb{Z}^n$ of prime power index}
\author{Hrishabh Mishra, Anwesh Ray}
\date{Source: arXiv:2211.16595v2 (CC BY 4.0)}
\begin{document}
\maketitle

\noindent\emph{Source and licence.} On the number of subrings of $\mathbb{Z}^n$ of prime power index. Version arXiv:2211.16595v2 (CC BY 4.0), distributed under the Creative Commons Attribution 4.0 International licence, as recorded in the arXiv licence metadata for this exact version and shown on the abstract page https://arxiv.org/abs/2211.16595v2. Full rights notice: licenses/arxiv-2211.16595-CC-BY-4.0.txt.\par\smallskip

\noindent\textit{Editorial scope.} Counted unit: the Lemma whose source label is \texttt{lemma 3} in the article (source0.tex lines 564--572, source label \texttt{lemma 3}), delivered together with the article's own complete original proof of it (source0.tex lines 574--643, one \texttt{proof} environment of 70 lines). No mathematics is written, repaired or re-derived: statement and proof are the article's own text line for line. Required context retained uncounted: lemma 2 (lines 518--522). Every definition, display and label of those lines is retained unchanged. Omitted from this package and not redistributed: 1--517, 523--563, 573--573, 644--1707. Nothing inside a retained excerpt is omitted or altered: every retained body is delivered exactly as the article's lines are written, so no omission receipt of any kind accompanies this package. Citations: the retained excerpts carry no citation command at all, so no bibliography entry of the article is transcribed and no reference is redistributed. Reference adaptation: none was needed and none was made; every reference inside the delivered unit resolves inside it, so no unresolved reference or citation is produced. Printed slips kept literal: no printed word, symbol, hypothesis or number of the counted statement or of its proof was altered. Layout and class: the article's own class is \texttt{amsart} with packages including geometry, caption, float, listings, todonotes, amscd, amsmath, amsxtra, amsthm, amssymb, stmaryrd, xr, mathrsfs, mathtools; the wrapper is the lane's pinned \texttt{amsart} editorial wrapper at 10pt on A4, which restates the theorem-like environments the retained text needs and adds the article's own macro definitions that the retained text needs (\texttt{\textbackslash op}), with extra wrapper packages (bm, bbm, dsfont, xspace, cancel, mathtools). The delivered numbering is the unit's own. Assets: no retained span contains a figure environment, a graphics-inclusion command, a PDF-page inclusion, a table or a TikZ picture; no image file is copied into this package, the package declares no asset dependency and no image file is redistributed with it. Licence: version arXiv:2211.16595v2 of this article is distributed under the Creative Commons Attribution 4.0 International licence, as recorded in the arXiv licence metadata for this exact version and shown on the abstract page https://arxiv.org/abs/2211.16595v2. A full rights notice is delivered at licenses/arxiv-2211.16595-CC-BY-4.0.txt. Characters: every retained excerpt is pure ASCII, so the delivered engine, XeLaTeX with its default Latin Modern text font, reports no missing character.
 \begin{lemma}\label{lemma 2}
 Let $\alpha$ be of the form $\alpha=(2,1,\dots, 1, 2, 1,\dots, 1, \beta)$, where $\beta>1$ and the second $2$ occurs at position $k\geq 2$, then
 \[g_\alpha(p)=p^{n-3+r}+p^{n-3}(p-1)r,\]
 where $r=n-1-k$.
 \end{lemma}

\begin{lemma}\label{lemma 3}
 Let $\alpha$ be of the form $\alpha=(2,1,\dots, 1, 2, 1,\dots, 1, \beta,1)$, where $\beta>1$, the second $2$ occurs at position $k\geq 2$ and set $r:=n-1-k$. Then, the following assertions hold.
 \begin{enumerate}
    \item If $\beta = 2$, then we find that
    \[g_\alpha(p)=2p^{n-3+r} + p^{n-5+r}(p-2) + 2rp^{n-3}(p-1) + p^{n-4}(p-1)(p-2)(p+r-2).\]
    \item If $\beta \geq 3$, then,
    \[g_\alpha(p)=2p^{n-3+r} + 2p^{n-5+r}(p-1) + 2rp^{n-3}(p-1) + 2p^{n-4}(p-1)^2(p+r-2).\]
\end{enumerate}
\end{lemma}

\begin{proof}
 Let $A$ be an integral matrix in Hermite normal form \[A=\begin{pmatrix}
	 p^{2}& p a_1 & pa_2 & \cdots & pa_{k-1}& \cdots & \cdots &\cdots& p a_{n-2} & 1 \\
	& p & 0 & \cdots & 0 &  \cdots & \cdots &\cdots& 0 & 1\\
	&& p & \cdots & 0 & \cdots &  \cdots &\cdots& 0  & 1\\
	&&& \ddots & \vdots& \vdots & \cdots &\cdots& \vdots & \vdots\\
	&&&& p^{2} & p b_1 & \cdots &\cdots& p b_r & 1\\
    &&&&  & p & \cdots &\cdots& 0 & 1\\
    	&&&&& & \ddots &\cdots& \vdots & 1\\
	&&&&&&& p^{\beta} & p c_1 & 1\\
	&&&&&&&&p&1\\
	&&&&&&&&&1.
	\end{pmatrix}\]
Since the first entry is $p^2$, it is clear that $v_i^2$ is contained in $\op{Col}(A)$ for $i=1,2,\dots,k$. For $i = 1,2,\dots,r-2$, we find that $v_{k+i}^2$ is contained in $\op{Col}(A)$ if and only if $v_{k+i}^2-pv_{k+i}$ is contained in $\op{Col}(A)$. So,
\[(a_{k-1+i}^2p^2-a_{k-1}p^2,\dots,b_i^2p^2-b_ip^2,\dots)^t\]
should be in $\op{Col}(A)$. It is clear that this is true if and only if
\begin{equation}
    a_{k-1}(b_i^2-b_i) \equiv 0 \mod{p}
\end{equation}
Similarly we see that $v_{n-2}^2 \in \op{Col}(A)$ if and only if $v_{n-2}^2 - p^{\beta}v_{n-2} \in \op{Col}(A)$. Now,
\[(a_{n-3}^2p^2-a_{n-3}p^{\beta+1},\dots,b_{r-1}^2p^2-b_{r-1}p^{\beta+1},\dots)^t\]
should be in $\op{Col}(A)$. It is again clear that this is true if and only if
\begin{equation}
    a_{k-1}b_{r-1} \equiv 0 \mod p
\end{equation}
as $\beta > 1$. Again, $v_{n-1}^2 \in \op{Col}(A)$ if and only if $v_{n-1}^2 - pv_{n-1} \in \op{Col}(A)$, so
\[(a_{n-2}^2p^2-a_{n-2}p^2,\dots,(b_r^2-b_r)p^2,\dots,(c_1^2-c_1)p^2, 0, 0)^t\]
should be in $\op{Col}(A)$. Then we have $c_1^2-c_1 \equiv 0 \mod p^{\beta-2}$, and
\begin{equation}
 b_{r-1}(c_1^2-c_1) \equiv 0 \mod p^{\beta-1}
 \end{equation}
 \begin{equation}
 a_{k-1}\left(b_r^2-b_r-\frac{b_{r-1}(c_1^2-c_1)}{p^{\beta-1}}\right)+a_{n-3}\frac{c_1^2-c_1}{p^{\beta-2}} \equiv 0 \mod p
\end{equation}
The requirement that $v_i\circ v_j$ belongs to $\op{Col}(A)$ for $i\neq j$ translates the following conditions on entries of the matrix (for all $1 \leq i < j \leq r$, and $\{i,j\} \neq \{r-1,r\}$)
\begin{equation*}
    a_{k-1}b_ib_j \equiv 0 \mod p
\end{equation*}
First, we suppose that $\beta \geq 3$, then we have conditions,
\begin{enumerate}
    \item $a_{k-1}(b_i^2-b_i) \equiv 0 \mod{p}$, $i=1,\dots,r-2$,
    \item $b_{r-1}a_{k-1} \equiv 0 \mod p$,
    \item $c_1^2-c_1 \equiv 0 \mod p^{\beta-2}$,
    \item $b_{r-1}(c_1^2-c_1) \equiv 0 \mod p^{\beta-1}$,
    \item $a_{k-1}\left(b_r^2-b_r-\frac{b_{r-1}(c_1^2-c_1)}{p^{\beta-1}}\right)+a_{n-3}\frac{c_1^2-c_1}{p^{\beta-2}} \equiv 0 \mod p$,
    \item $a_{k-1}b_ib_j \equiv 0 \mod p$, $1 \leq i < j \leq r$, and $\{i,j\} \neq \{r-1,r\}$.
\end{enumerate}
We consider two cases.
\par \textbf{Case 1 :} First, we consider the case when $a_{k-1} = 0$. Then, the above equations reduce to the following:
\begin{enumerate}
    \item $c_1^2-c_1 \equiv 0 \mod p^{\beta-2}$,
    \item $b_{r-1}(c_1^2-c_1) \equiv 0 \mod p^{\beta-1}$,
    \item $\frac{a_{n-3}(c_1^2-c_1)}{p^{\beta-2}} \equiv 0 \mod p$.
\end{enumerate}
If $c_1 \in \{0,1\}$, then the number of such matrices is $2p^{n-3+r}$. On the other hand, if $c_1 \notin \{0,1\}$, then, there are $2p-2$ choices for $c_1$. Since $b_{r-1} = 0, a_{n-3}=0$, we deduce that there are $2(p-1)p^{n+r-5}$ more matrices.
\par \textbf{Case 2 :} Suppose $a_{k-1} \neq 0$, then equations reduce to
\begin{enumerate}
\item $b_i^2-b_i \equiv 0 \mod p$, $i =1, \dots, r-2$,
\item $b_{r-1}=0$,
\item $c_1^2-c_1 \equiv 0 \mod p^{\beta-2}$,
\item $a_{k-1}\left(b_r^2-b_r\right)+a_{n-3}\frac{c_1^2-c_1}{p^{\beta-2}} \equiv 0 \mod p$,
\item $b_ib_j \equiv 0 \mod p$, $1 \leq i < j \leq r$, and $\{i,j\} \neq \{r-1,r\}$
\end{enumerate}
We divide into two sub-cases.
\begin{itemize}
    \item \textbf{Case 2A :} Suppose that $c_1 \in \{0,1\}$, then we get that $b_r \in \{0,1\}$ and as in Lemma \ref{lemma 2}, we get that at most one of the $b_i$ satisfies $b_i=1$. Subdivide into two cases, first consider the case when all the $b_i$ are equal to $0$ and then consider the case when at most one of the $b_i$ is equal to $1$. Therefore, the total number of choices is $2rp^{n-3}(p-1)$.
    \item \textbf{Case 2B :} Consider the other subcase, i.e., $c_1 \notin \{0,1\}$. Then, there is a unique value of $a_{n-3}$ for each value of $b_r, a_{k-1}, c_1$, so we get $2p^{n-3}(p-1)^2(p+r-2)$ matrices
\end{itemize}
Putting it all together, the result is proven. The case $\beta = 2$ is similar, and the number of matrices change only in Case 1, with $c_1\notin \{0,1\}$ and Case 2B.
\end{proof}

\end{document}
